% Industry resume, Grover Aruquipa
% Compile with pdfLaTeX (also works on Overleaf). One column on purpose:
% applicant tracking systems read single-column text far more reliably than sidebars.
\documentclass[10pt,letterpaper]{article}
\usepackage[T1]{fontenc}
\usepackage[utf8]{inputenc}
\usepackage[scaled=0.95]{helvet}
\renewcommand{\familydefault}{\sfdefault}
\usepackage[top=1.0cm,bottom=0.8cm,left=1.5cm,right=1.5cm]{geometry}
\usepackage{enumitem}
\usepackage{titlesec}
\usepackage{xcolor}
\usepackage{microtype}
\usepackage[hidelinks]{hyperref}

\definecolor{ink}{HTML}{111111}
\definecolor{accent}{HTML}{0D2B61}
\definecolor{muted}{HTML}{555555}
\color{ink}
\pagestyle{empty}
\setlength{\parindent}{0pt}
\setlength{\parskip}{0pt}

\titleformat{\section}{\normalsize\bfseries\color{accent}}{}{0em}{\MakeUppercase}[\vspace{-6pt}{\color{accent}\rule{\linewidth}{0.6pt}}]
\titlespacing*{\section}{0pt}{5pt}{2pt}
\setlist[itemize]{leftmargin=1.1em,itemsep=0.5pt,topsep=1pt,parsep=0pt,label=\textbullet}

% \role{Title}{Organization}{Location}{Dates}
\newcommand{\role}[4]{\textbf{#1}, #2 \hfill {\color{muted}#3 \textbar\ #4}\par}

\begin{document}

%================ HEADER ================
{\LARGE\bfseries Grover Aruquipa}\par\vspace{2pt}
{\large Mechanical Engineer, Robotic Surgical Accessories \textbar\ Precision Mechanisms, Concept to Verified Design}\par\vspace{3pt}
{\small Qu\'ebec, QC, Canada \ \textbullet\ \
\href{mailto:grover.aruquipa.9@gmail.com}{grover.aruquipa.9@gmail.com} \ \textbullet\ \
\href{https://www.linkedin.com/in/grover-aruquipa-996917126}{linkedin.com/in/grover-aruquipa-996917126} \ \textbullet\ \
\href{https://groveraruquipa.github.io}{groveraruquipa.github.io}}\par\vspace{1pt}
{\small\bfseries Available from January 2027 \ \textbullet\ \ Open to relocation}\par


%================ SUMMARY ================
\section{Summary}
Mechanical engineer (MSc, PhD defense early 2027, Universit\'e Laval) who designs, builds and verifies precision robotic
mechanisms. I have taken three robots from requirements and architecture through SolidWorks design, stress analysis,
component selection, fabrication, assembly and verification testing, and I am building a fourth. The most compact is a
4-DoF coaxial wrist-gripper that rolls its tool without limit while the jaw is driven through concentric hollow shafts,
the same function a wristed surgical instrument performs.

%================ SKILLS ================
\section{Technical Skills}
\begin{itemize}
\item \textbf{Mechanism design:} SolidWorks (expert), Creo; precision joints, coaxial and concentric transmissions,
bearings, integrated grippers; component selection (motors, encoders, gearboxes, bearings); design for manufacturability,
assembly and reliability.
\item \textbf{Analysis:} statics, dynamics, kinematics, strength of materials and stress analysis; FEA in ANSYS; Jacobians,
singularity and workspace analysis; MATLAB, Simscape, MuJoCo, NVIDIA Isaac Sim, Pinocchio.
\item \textbf{Prototyping and verification:} machine shop work, CNC machining and 3D printing (SLS, FDM); fixtures and test
setups; component verification and system validation tests; measurement of accuracy and repeatability.
\item \textbf{Documentation and software:} design specifications, test reports and peer-reviewed technical documentation;
MS Office; real-time control at 500\,Hz (C++, Simulink Real-Time), Python, ROS~2.
\end{itemize}

%================ EXPERIENCE ================
\section{Experience}
\role{Doctoral Researcher}{Robotics Laboratory, Universit\'e Laval}{Qu\'ebec, Canada}{2023 -- present}
\begin{itemize}
\item \textbf{Led the full development of three robotic mechanisms} (4, 7 and 8 actuators): requirements, architecture,
detailed design, analysis, fabrication, assembly, verification testing and redesign.
\item \textbf{4-DoF coaxial wrist-gripper:} four actuation channels through concentric hollow shafts, non-backdrivable
jaw, 180\,mm footprint, 0.25\,kg moving mass; verified over ten continuous revolutions with sub-degree tracking.
\item \textbf{(6+1)-DoF parallel robot:} gripper actuated from the base through a concentric tube while the end effector
rotates without limit. \textbf{Delta robot with Sarrus gripper (in progress):} sensorless grasp-force control.
\item Presented each design for review before fabrication; mentored two graduate students day to day on an 8-actuator
robot; first author of 3 journal papers and 4 under review; reviewer for ASME J.\ Mechanisms and Robotics.
\end{itemize}

\vspace{3pt}
\role{Visiting Researcher}{Institute of Robotics, Johannes Kepler University Linz}{Linz, Austria}{Apr. -- May 2026}
\begin{itemize}
\item Developed singularity and grasp-quality indices for manipulators with integrated grasping, hosted by Univ.-Prof.~A.~M\"uller (Mitacs Globalink Research Award).
\end{itemize}

\vspace{3pt}
\role{Research Intern}{FEMTO-ST (CNRS)}{Besan\c{c}on, France}{2022 -- 2023}
\begin{itemize}
\item Designed and modeled a redundant dexterous parallel robot with two unlimited rotations for micro-manipulation, and wrote its real-time C++ control software; clean-room fabrication projects.
\end{itemize}

\vspace{3pt}
\textbf{Industry roles} \hfill {\color{muted}2020 -- 2021}\par
{\small Computer Vision Engineer, Sencon Srl (industrial vision, scaffolding detection) \ \textbullet\ \ Data Science
Intern, Hitachi ABB Power Grids, Zurich (machine learning for high-voltage circuit breakers) \ \textbullet\ \ Data
Scientist, KTH Royal Institute of Technology.}

%================ EDUCATION ================
\section{Education}
\textbf{PhD, Mechanical Engineering}, Universit\'e Laval \hfill {\color{muted}Qu\'ebec, Canada \textbar\ 2023 -- 2027 (expected)}\par
{\small Kinematically redundant parallel robots with unlimited rotations and integrated grasping. Advisors: C.~Gosselin, P.~Lambert.}\par\vspace{2pt}
\textbf{MSc, Smart Integrated Systems}, Universit\'e Bourgogne--Franche-Comt\'e \hfill {\color{muted}Besan\c{c}on, France \textbar\ 2021 -- 2023}\par\vspace{2pt}
\textbf{BEng, Mechatronics Engineering}, Bolivian Catholic University \hfill {\color{muted}La Paz, Bolivia \textbar\ 2014 -- 2019}\par
{\small Thesis: reconfigurable Delta robot for product picking.}

%================ AWARDS ================
\section{Awards and Languages}
\begin{itemize}
\item 1st Place, Graduate Division, ASME Student Mechanisms and Robotics Design Competition (2025), team lead.
\item Finalist, MassRobotics competition, Boston, with one of my robot prototypes, team lead.
\item Mitacs Globalink Research Award (2026); 1st Place, Student Paper Competition, Argentine Robotics Conference (2019).
\item Languages: Spanish (native), English (full professional), French (intermediate).
\end{itemize}

\newpage
%================ PUBLICATIONS
\section{Selected Publications}
{\small First author on all papers. Full list at groveraruquipa.github.io.}\par\vspace{2pt}
\begin{itemize}
\item \textbf{G.~Aruquipa}, P.~Lambert, C.~Gosselin, ``A compact coaxial parallel wrist-gripper with unlimited tool-axis rotation,'' under review, \emph{Robotics and Autonomous Systems}.
\item \textbf{G.~Aruquipa}, P.~Lambert, C.~Gosselin, ``Analysis and design of a novel (6+1)-DoF parallel robot with unlimited rotation and grasping capabilities,'' \emph{Mechanism and Machine Theory}, vol.~224, art.~106438, 2026.
\item \textbf{G.~Aruquipa}, V.~Muralidharan, P.~Lambert, C.~Gosselin, ``Analysis and design of a novel compact four-degree-of-freedom parallel robot with unlimited rotation,'' \emph{ASME Journal of Mechanisms and Robotics}, vol.~18, no.~6, art.~061001, 2026.
\item \textbf{G.~Aruquipa}, P.~Lambert, C.~Gosselin, ``An analytical framework for collision-free and singularity-free unlimited rotation in spatial parallel mechanisms,'' under review, \emph{Mechanism and Machine Theory}.
\item \textbf{G.~Aruquipa}, P.~Lambert, R.~Dahmouche, C.~Gosselin, ``Parallel robots for grasping tasks: a review,'' \emph{ASME Journal of Mechanisms and Robotics}, vol.~18, no.~12, art.~120801, 2026. doi:10.1115/1.4072612
\end{itemize}

%================ REFERENCES ================
\section{References}
\begin{itemize}
\item \textbf{Prof.~Cl\'ement Gosselin}, Full Professor, Department of Mechanical Engineering, Universit\'e Laval, Qu\'ebec,
Canada. PhD supervisor. \href{mailto:clement.gosselin@gmc.ulaval.ca}{clement.gosselin@gmc.ulaval.ca}
\item \textbf{Prof.~Patrice Lambert}, Assistant Professor, Department of Mechanical Engineering, Universit\'e Laval,
Qu\'ebec, Canada. PhD co-supervisor. \href{mailto:patrice.lambert@gmc.ulaval.ca}{patrice.lambert@gmc.ulaval.ca}
\item \textbf{Univ.-Prof.~Dr.-Ing.~Andreas M\"uller}, Head of the Institute of Robotics, Johannes Kepler University Linz,
Austria. Host during my research stay (2026). \href{mailto:a.mueller@jku.at}{a.mueller@jku.at}
\end{itemize}

\end{document}
